\documentclass[tikz,border=6pt]{standalone}
\usepackage{ctex}
\usepackage{amsmath}
\usepackage{amssymb}
\usepackage{pgfplots}
\pgfplotsset{compat=1.18}
\usetikzlibrary{calc,angles,quotes,arrows.meta,positioning,decorations.markings}
\begin{document}
\begin{tikzpicture}
  \begin{axis}[
      width=11.5cm, height=11.5cm,
      axis equal image,
      xmin=-2.4, xmax=4.6, ymin=-2.4, ymax=4.6,
      axis lines=middle,
      xlabel={$x$}, ylabel={$y$},
      xtick={-2,-1,1,2,3,4}, ytick={-2,-1,1,2,3,4},
      ticklabel style={font=\footnotesize},
      clip=false,
      every axis x label/.style={at={(rel axis cs:1,0.5)},anchor=west},
      every axis y label/.style={at={(rel axis cs:0.5,1)},anchor=south},
    ]

    % 对称轴 y = x
    \addplot[domain=-2.2:4.4,samples=2,gray!60,dashed,thick] {x};
    \node[gray!65!black,font=\small] at (axis cs:3.9,4.2) {$y=x$};

    % 指数函数 a^x （取 a=e 作代表，a>1）
    \addplot[domain=-2.3:1.5,samples=120,thick,blue!70!black] {exp(x)};
    \node[blue!70!black] at (axis cs:1.62,4.2) {$y=a^{x}$};

    % 对数函数 log_a x （与上面关于 y=x 对称）
    \addplot[domain=0.12:4.4,samples=160,thick,red!70!black] {ln(x)};
    \node[red!70!black] at (axis cs:4.25,1.7) {$y=\log_a x$};

    % e 的标注：指数曲线过 (1,e)，对数曲线过 (e,1)
    \filldraw[blue!70!black] (axis cs:1,2.718) circle (1.4pt);
    \node[blue!60!black,font=\footnotesize,left] at (axis cs:1,2.718) {$(1,e)$};
    \filldraw[red!70!black] (axis cs:2.718,1) circle (1.4pt);
    \node[red!60!black,font=\footnotesize,below right=-1pt] at (axis cs:2.718,1) {$(e,1)$};

    % 一对对称点示意：(0,1) 与 (1,0)
    \filldraw[blue!70!black] (axis cs:0,1) circle (1.4pt);
    \filldraw[red!70!black] (axis cs:1,0) circle (1.4pt);
    \draw[teal!55!black,thin,dash pattern=on 2pt off 2pt] (axis cs:0,1) -- (axis cs:1,0);

    % 半透明注释框：综合主线交汇（放右下空白区，避开曲线）
    \node[fill=teal!10,draw=teal!55!black,rounded corners,align=left,
          font=\footnotesize,inner sep=6pt,fill opacity=0.85,text opacity=1]
      at (axis cs:2.95,-1.5)
      {\textbf{五条主线交汇}\\[2pt]
       指数 $\leftrightarrow$ 对数（反函数）\\
       自然底数 $e$ 连接微积分\\
       可加化 $\cdot$ 尺度压缩\\
       信息熵 $\cdot$ 复对数};

    % 关于 y=x 互为镜像的说明（左上空白）
    \node[fill=gray!10,draw=gray!50,rounded corners,align=center,
          font=\footnotesize,inner sep=4pt,fill opacity=0.85,text opacity=1]
      at (axis cs:-0.9,3.3)
      {$a^{x}$ 与 $\log_a x$\\ 关于 $y=x$ 对称};

  \end{axis}
\end{tikzpicture}
\end{document}
